\begin{center}
Problema A

Abóbora Perfeita

{\small Nome do arquivo: A.c, A.cpp, A.java, ou A.py}

{\large Timelimit: $1$}

{\tiny Por Athos José}

\end{center}

Na cidade de Rio Verde, a produção de abóboras é tão abundante que, ao longo dos anos, surgiram várias lendas sobre abóboras especiais. Uma das lendas mais antigas fala sobre a mítica Abóbora Perfeita.

Segundo a lenda, uma abóbora só pode ser considerada perfeita quando nascem, ao redor da abóbora maior, abóboras menores cujos respectivos tamanhos sejam divisores da abóbora maior. Além disso, a soma dos tamanhos das abóboras menores ao redor deve resultar no tamanho da abóbora maior. Essas propriedades tornam a abóbora perfeita.

Você foi chamado para investigar abóboras de diferentes tamanhos e determinar quais delas podem ser consideradas perfeitas, de acordo com a lenda.

\vspace{0.3cm}
\textbf{Entrada}

A primeira linha contem um inteiro $N$, representando o número de abóboras $(1 \leq N \leq 10^7)$. As $N$ linhas seguintes contem um inteiro $T$ $(1 \leq T \leq 10^4)$, onde T indica o tamanho da abóbora.

\vspace{0.3cm}
\textbf{Saída}

Para cada caso de teste imprima ``SIM'' se a abóbora pode ser considerada perfeita. Caso contrario imprima ``NAO''.

\begin{table}[ht]
    \centering
    \begin{tabular}{|p{7cm}|p{7cm}|}
        \hline
        \begin{center}
            \vspace{-0.3cm}
            \textbf{Exemplos de Entrada}
            \vspace{-0.3cm}
        \end{center} 
         &  
        \begin{center}
            \vspace{-0.3cm}
            \textbf{Exemplos de Saída}
            \vspace{-0.3cm}
        \end{center} \\ \hline
         \texttt{4} & \texttt{NAO} \\
         \texttt{9} & \texttt{SIM} \\
         \texttt{6} & \texttt{SIM} \\
         \texttt{28} & \texttt{NAO} \\
         \texttt{21} & \texttt{} \\
         \hline
    \end{tabular}
    
    \caption{Exemplos de entradas e saídas}
    \label{tab:inputA3}
    
\end{table}


\vspace{0.3cm}
\textbf{Nota}

O exemplo contem $4$ caso de testes.

No \textbf{primeiro} caso temos $1+3 = 4$ que é diferente do tamanho $9$.

Já no \textbf{segundo} caso temos $1+2+3 = 6$ que é igual ao tamanho da abóbora $6$.

\newpage

\begin{center}
\noindent
\lstinputlisting{abobora_24/A/A4.cpp}
\end{center}